% Ч1.4 — Увеличение и уменьшение на p% (множитель 1 ± p/100)
% Старая методичка: стр. 152 (старое задание 16)
% Источники: методичка по заданию 13 (часть 1). Стиль — стиль-изложения.md

\begin{tema}{Ч1.4}
Чтобы \opr{увеличить} число на $p\%$, можно умножить его на $1+\frac{p}{100}$:
\[
  x+\frac{p}{100}\cdot x=\left(1+\frac{p}{100}\right)x.
\]
Этот подход работает и для \opr{уменьшения}: чтобы уменьшить число на $p\%$,
нужно умножить его на $1-\frac{p}{100}$. Например, $+23\%$ — умножить на $1{,}23$,
$-45\%$ — на $0{,}55$.

Число $k=1\pm\frac{p}{100}$ называют \opr{повышающим} (\opr{понижающим}) коэффициентом.
Если число меняется несколько раз, коэффициенты перемножаются.

\begin{primer}
Холодильник стоит 15\,000 рублей. Сколько он будет стоить после увеличения
цены на 27\,\%? Решение: $1{,}27\cdot15\,000=19\,050$. Ответ: 19\,050 рублей.
\end{primer}

На этом построена экономическая задача (№\,13): каждый год долг или вклад
умножается на $k$ — см. \tref{М2.1}.

\begin{oshibka}
Цену подняли на 20\,\%, а потом снизили на 20\,\% — она \emph{не} вернулась
к исходной: $1{,}2\cdot0{,}8=0{,}96$, то есть цена упала на 4\,\%.
\end{oshibka}
\end{tema}
